% ============================================================
%  Harald Høffding, « Pascal et Kierkegaard »
%  Revue de Métaphysique et de Morale, 30e année, no 2
%  (avril--juin 1923), pp. 221--246.
%  The Pascal tercentenary number ("Troisième centenaire de la
%  naissance de Pascal").
%  TRANSCRIPTION (French --- Høffding wrote this article in French)
%  Source: Internet Archive item revue-de-metaphysique-et-de-morale-30-2,
%          Public Domain Mark 1.0. Local copy: scan.pdf (gitignored).
%          https://archive.org/details/revue-de-metaphysique-et-de-morale-30-2
%  PAGE MAP:  see pagemap.py.  OFFSET NOT YET VERIFIED.
% ============================================================
%
%  PAGE-JOINT AND PARAGRAPH REPAIR (2026 pass; fifth item swept, after
%        speculative-methode where the fault was found, evangelietroen-
%        theologien, om-personlig-sandhed and johannesclimacos):
%        Both files here were built by splicing page-range fragments into a
%        skeleton that separated them with blank lines. A blank line is a
%        paragraph break, so wherever the printed text ran on across a leaf the
%        splice invented a paragraph the article does not have.
%        Repaired in this file:
%          * 2 page-joint hyphens -> `\-%` (pp. 232 "con-"/"sidérant" and 236
%            "concep-"/"tion")
%          * 9 FALSE paragraph breaks removed (pp. 231--239, nine consecutive
%            leaves), every one of them plainly mid-sentence
%          * 0 missing paragraph breaks
%          * SEPARATELY, 8 printed LINE-break hyphens that had been left bare
%            mid-page and printed as "précipita- tion", "pour- quoi",
%            "primi- tive", "mar- quée", "souf- france", "pro- vincial",
%            "chris- tianisme", "pro- blèmes". These are not the page-joint
%            fault at all — they are ordinary line divisions in the printed
%            French that were never closed up — but they produce exactly the
%            same artefact and are fixed the same way. The `\emph{}` spanning
%            the "pro-"/"vincial" break was preserved.
%        `paragraphs.py` in this directory performs the audit and takes either
%        file as its argument. Run it before any future publication.
%        ITS ONE HARD PROBLEM HERE IS THE RUNNING HEAD, which is full-measure
%        and lies in BOTH directions: on even pages the Revue sets the folio out
%        in the left margin, so the head begins ~340 px LEFT of the block and the
%        page reads as a -337 px hanging indent; on odd pages the head is
%        CENTRED, so it begins ~100 px RIGHT of the block and is
%        indistinguishable from a paragraph indent — printed p. 245 was reported
%        as +97 px on that mistake and in fact opens flush. Two guards are needed
%        together: the line must sit in the text block AND not be followed by the
%        extra white space that separates a head from the body. With both,
%        continuations measure -19 to +20 px and the one paragraph opening +77.
% ============================================================
\documentclass[12pt,a4paper]{article}
\usepackage[utf8]{inputenc}
\usepackage[T1]{fontenc}
\usepackage{libertinus}
\usepackage{libertinust1math}
\usepackage{textalpha}   % Greek typed directly, if Høffding quotes any
\usepackage[danish,main=french]{babel}   % Danish for the Kierkegaard quotations
\usepackage{csquotes}
\usepackage[protrusion=true,expansion=true]{microtype}
\usepackage{setspace}
\usepackage[margin=1.2in]{geometry}
\usepackage{enumitem}
\usepackage{fancyhdr}
\setlength{\headheight}{28pt}
\addtolength{\topmargin}{-13.5pt}

% Danish inside the French text. Use \dk{...} so the hyphenation and any
% babel-sensitive punctuation follow the quoted language, not the frame.
\newcommand{\dk}[1]{\foreignlanguage{danish}{#1}}

\usepackage{hyperref}
\hypersetup{
  pdftitle={Pascal et Kierkegaard},
  pdfauthor={Harald Høffding},
  hidelinks
}

\pagestyle{fancy}
\fancyhf{}
\fancyhead[LE,RO]{\thepage}
\fancyhead[RE]{\textit{Revue de Métaphysique et de Morale}}
\fancyhead[LO]{\textit{Pascal et Kierkegaard}}

\setstretch{1.3}

\begin{document}

\begin{center}
  {\Large PASCAL ET KIERKEGAARD}\\[1ex]
  % [verify the byline and any standfirst against the image of p. 221]
\end{center}

\bigskip

% --- The article is in four parts. Their boundaries are NOT detectable from
% --- the OCR (no standalone roman numerals survive it), so the skeleton does
% --- not guess at them: each batch agent sets the section heads it actually
% --- sees on the page, inline in its own fragment. The three markers below are
% --- equal thirds for dispatch only and carry no structural claim.

% --- p. 221 ---
%
% TITLE BLOCK ON p. 221, as actually printed (read off the image, PDF 239):
%   - a DOUBLE RULE across the measure at the head of the page;
%   - the title, centred, in letterspaced full capitals: PASCAL ET KIERKEGAARD
%     (no italic, no small caps, no accents involved);
%   - a short centred rule beneath it;
%   - then the text.
% There is NO byline on p. 221, NO standfirst, NO epigraph and NO footnote
% hung off the title. Nothing announces the Pascal tercentenary on this page.
% The author's name appears only in the recto running head, set
%   H. HÖFFDING. --- PASCAL ET KIERKEGAARD.   <page no.>
% with "H. HÖFFDING." in bold full caps (umlaut Ö) and the rest in small caps;
% verso heads read   <page no.>   REVUE DE MÉTAPHYSIQUE ET DE MORALE.
% The title block itself is set by the skeleton above this marker; it is only
% recorded here, not repeated.

A deux siècles de distance, deux penseurs, également attachés
pourtant au christianisme jusqu'au tréfond de leurs âmes ferventes,
se sont rencontrés dans le procès de leur religion telle que les temps
l'ont déformée. L'un et l'autre dénoncent l'incompatibilité qui
oppose l'esprit qui animait le christianisme à ses débuts et son état
des temps modernes. Cet antagonisme date du jour où l'Église a
tenté d'entrer en relations intimes avec le « monde », le « siècle »,
autrement dit avec la culture développée sur des bases purement
humaines. Le bilan de l'Église fut entrepris, dans les deux cas, avec
tant de passion, et poussé avec tant de logique qu'il y aurait eu
matière à une grande révolution intellectuelle et morale, mais ---
et c'est bien là le point tragique des destinées en cause --- à peine
les deux penseurs eurent-ils énoncé leur dernière parole, qu'ils
succombèrent, consumés l'un et l'autre par la lutte spirituelle qu'ils
avaient eue à soutenir avec eux-mêmes et avec leurs contemporains.
Ils moururent jeunes, Pascal à 39 ans (1662) et Kierkegaard à 42
(1855), laissant entier le grand problème formulé à nouveau par
eux avec une acuité sans précédent. Mais, dans le monde de l'esprit,
jamais l'énergie héroïque ne se dépense en vain et, reprise avec
cette instance et cette fougue, la question s'est posée de savoir si,
oui ou non, les traditions religieuses des Églises chrétiennes pouvaient
continuer d'appuyer l'inquiétude loyale et la recherche sincère
des amis de la vérité.

Vaine question ! du camp des Églises n'est venue aucune réponse.
% ^ printed thus: exclamation mark, then a LOWER-CASE "du". Not a defect
%   we may mend; transcribed as set.
On s'y flatte de persister dans la pratique suivie en négligeant les récriminations.
On s'appuie sur le fait, d'ailleurs incontestable, qu'il n'a
pas paru jusqu'à présent d'autre organisation susceptible à la fois de
répandre dans le monde la notion du sérieux de l'existence humaine
% --- p. 222 ---
et d'apporter aux hommes l'allégement et le réconfort que sa dureté
exige. Ainsi le problème reste en suspens. On raconte de Scipion
l'Africain qu'accusé devant les comices d'irrégularités dans ses
comptes, il s'écria : « Romains, à pareil jour, j'ai vaincu Annibal à
Zama ; montons au Capitole en rendre grâces aux dieux ! » Le peuple
le suivit, et il ne fut plus question de reddition de comptes. Le
doute a gardé toute sa force depuis les deux penseurs et les études
historiques en ont même appronfondi la base.
% ^ PRINTER'S ERROR, as printed: "appronfondi" for "approfondi".
%   Verified at 300 dpi and again on a 2x enlargement of the line.
Les Églises ont-elles
le droit de croire à une continuité vivante entre elles et le grand
mouvement spirituel, parti il y a plus de dix-neuf siècles d'un coin
de l'empire romain et dont l'idée maîtresse commandait de vivre
désormais dans l'attente extrême du proche et surnaturel avènement
d'un règne nouveau, attente qui rabaissait à une valeur infime
les cadres humains de l'existence, la famille et l'État, l'art et la
science, quand elle n'y voyait pas autant d'obstacles au royaume de
Dieu ?

Ce n'est d'ailleurs pas ce procès seulement des Églises, quelque
gravité qu'il ait eue, qui nous fait souvenir de Pascal et de Kierkegaard.
Avant de conclure, leurs esprits ont développé des idées de
grande importance : ici, comme souvent ailleurs dans le domaine de
la pensée, la route parcourue a sa pleine valeur, indépendamment
du but atteint. Chemin faisant, ils ont projeté des clartés sur la
nature et les conditions d'existence de la vie spirituelle de l'homme,
moissonnant sans compter des trouvailles psychologiques d'un haut
prix. Ils ont apporté des contributions à ce que nous appelons
aujourd'hui la théorie de la connaissance en faisant jouer la pensée
à la dernière limite de son domaine. Ils ont montré que là même
où le sentiment et la passion portent la parole, c'est encore l'idée
qui détermine --- telle une étoile conductrice --- la direction du
mouvement psychologique. Ils sont les champions de la Pensée, et
leur lutte reste intéressante, qu'elle ait abouti ou non à la victoire.
Ils suivirent des chemins divers. Leurs individualités apparaissent,
% ^ "Ils suivirent des chemins divers." is NOT a new paragraph: the line is
%   set flush left, no indent. Same paragraph, on the image.
au dedans et au dehors, sous des aspects différents, présentant,
chacune, des traits constitutifs qui lui sont propres, et les milieux
historiques où s'aiguisèrent leurs pensées étaient également fort
dissemblables, tout en offrant entre eux certaines analogies.

Par leur style, déjà, ils diffèrent. Nous devons nous rappeler ici
que les \emph{Pensées} de Pascal, qui constituent son œuvre principale,
ne se présentent pas à nous sous la forme définitive d'un ouvrage
% --- p. 223 ---
mené à bonne fin par son auteur. Ce sont des fragments, des notes
griffonnées par Pascal dans les répits que lui laissait sa maladie,
interrompues souvent par un début de crise. Concis, resserrés,
vibrant d'une émotion qui, dans un rythme alternatif, fait jaillir les
pensées ou jaillit d'elles, ces fragments ont une force et une clarté
sans égale, et il n'y a pas lieu de croire que ces qualités se fussent
perdues dans une rédaction définitive. Le style de Kierkegaard n'a
pas cette frappe puissante ni cette densité. Pour lui --- du moins
dans la première phase de sa production littéraire --- le travail de
composition était une cure : tant qu'il rédigeait, la mélancolie
n'avait pas de prise sur son âme. Aussi prenait-il son temps. Et
puis, il aimait la langue pour elle-même, il jouait avec elle sans
prendre garde, quelquefois, que le lecteur, lui, était peut-être un
peu plus pressé, impatient d'embrasser dans son ensemble la suite
des idées, de voir à quoi tout cela pourrait bien aboutir. Des digressions
lyriques, des descriptions de la nature ou des épisodes empruntés
à la vie de tous les jours viennent s'entrelacer à la chaîne des
développements. La passion, qui est à la base de sa production,
comme de celle de Pascal, se fait sentir plutôt sous la forme d'un
frémissement continu que sous celle de crises émotives. Ce n'est
que dans la dernière période de sa production que ce frémissement
se change en tremblement de terre.

Si, dans cette étude comparative, il m'arrive de parler un peu plus
longuement du penseur danois, c'est que je m'adresse à des lecteurs
français. En France, l'étude de Pascal a suscité une production tellement
abondante, tellement solide et documentée que prétendre
fournir un apport ce serait vraiment offrir des chouettes aux Athéniens.
Qu'il me soit permis, à ce propos, d'exprimer mon admiration
pour le \emph{Port-Royal} de Sainte-Beuve, que je viens de relire à
près de 60 ans de distance et qui, après tant de travaux parus depuis,
garde toute sa fraîcheur et toute sa pénétration. Il faut bien aussi
que je rende hommage à cette splendide édition des \emph{Œuvres} de Pascal
que nous devons à M. Léon Brunschvicg et à ses collaborateurs
MM. Pierre Boutroux et Félix Gazier et dont l'\emph{Introduction} et les
annotations fournissent des éclaircissements d'une grande valeur.
Je sais gré surtout à M. Brunschvicg d'avoir, dans une \emph{Introduction}
aux \emph{Pensées}, exposé de façon extrêmement intéressante la suite
logique des idées qui aurait probablement été celle de Pascal dans
l'\emph{Apologie} qu'il se proposait d'écrire. Dans les lignes qui suivent
% --- p. 224 ---
je m'en rapporterai, pour les citations des \emph{Pensées}, à l'ordre proposé
par M. Brunschvicg. --- Quant aux ouvrages de Kierkegaard,
une seconde édition de ses œuvres complètes est actuellement en
cours de publication. Il en a paru une édition allemande (chez Diederichs,
Iéna). Elles comprennent quatorze volumes. Mais, en outre,
on a commencé la publication intégrale du \emph{Journal} et des papiers
intimes de Kierkegaard (par P.-A. Heiberg et V. Kuhr). Jusqu'à
présent onze volumes ont paru et il reste encore la matière de sept
volumes. Nous sommes donc en présence d'une production très volumineuse.
En français, une belle étude sur Kierkegaard par M. Delacroix
a été publiée, en 1900, dans la présente \emph{Revue}, qui a donné
aussi, en 1913, un discours prononcé par l'auteur de ces lignes à
l'Université de Copenhague pour célébrer le centenaire du philosophe
danois (5 mai 1913). Une monographie sur Kierkegaard
philosophe, que j'ai écrite en 1892, a été traduite en allemant et va
% ^ PRINTER'S ERROR, as printed: "allemant" for "allemand".
être traduite en espagnol.

L'hypothèse d'une influence exercée par Pascal sur Kierkegaard
paraît moins que probable. Il est vrai que les œuvres de Pascal,
en traduction allemande, figurent dans l'inventaire de la bibliothèque
de Kierkegaard, mais ce n'est que sur le tard (en automne et
en hiver 1850\footnote{D'après un renseignement que je dois à l'obligeance de M. P.-A. Heiberg,
l'un des éditeurs du \emph{Journal}.}) que ce dernier fait mention de Pascal dans son
% ^ FOOTNOTE-MARK SCHEME, first note of the article, on p. 224:
%   in the text the mark is a superscript ARABIC numeral immediately followed
%   by a CLOSING PARENTHESIS on the baseline --- the page reads "1850 ¹)".
%   At the foot the note is set in small type, indented, opening "1. D'après",
%   i.e. numeral + full stop, and there is NO rule above it (blank space only).
%   Only one note falls in pp. 221--229, so whether the numbering restarts per
%   page or runs on cannot be settled from this batch.
\emph{Journal}, --- c'est-à-dire à un moment où il avait déjà donné les
écrits [\emph{La Maladie à la mort} (1849) et \emph{L'Exercice dans le Christianisme}
(1850)] qui présentent le plus d'affinité avec les idées de
Pascal. Il a lu, depuis, divers ouvrages traitant de Pascal, notamment
celui de Reuchlin, relatif au Port-Royal, et il a remarqué la
différence qui existe entre les anciennes éditions des \emph{Pensées} et
celle due à Faugère où se trouvent admises les notes les plus révolutionnaires.
(Les traductions allemandes représentées dans la
bibliothèque de Kierkegaard avaient toutes été faites sur les premières
éditions.)

Dans le rapprochement des deux penseurs que je vais entreprendre,
je m'arrêterai successivement à leur tempérament et à leur
caractère, à leurs antécédents intellectuels, à leur rapport au christianisme
et à la question de savoir quelles voies ils laissaient
ouvertes au moment où ils se turent.
% --- p. 225 ---
% Section head, centred, at the top of p. 225 immediately under the running
% head. Set as:  I. --- Leur tempérament et leur caractère.
% with the numeral and dash in full caps, the first letter of each word in
% full caps and the rest in SMALL CAPS, and a full stop at the end.
% This is the article's FIRST section head: pp. 221--224 are an unnumbered
% introduction, and section I opens here.
\section*{I. --- \textsc{Leur tempérament et leur caractère.}}

Le tempérament physique et moral de Pascal a été défini par ses
deux sœurs de façon concordante (lettres de Jacqueline, biographie
de M\textsuperscript{me} Perier). Il était d'humeur bouillante et la vivacité de son
esprit n'allait pas sans une impatience qui le rendait, par moments,
difficile à satisfaire ; sa sœur aînée admirait qu'avec la fougue de
son tempérament il eût été préservé des écarts de la jeunesse. Une
pensée qui me paraît tout à fait significative à cet égard est celle
du fragment 139 (Br.) selon laquelle tout le malheur des hommes
viendrait d'une seule chose, qui est de « ne savoir pas demeurer
en repos ». Dans le curieux \emph{Discours sur les passions de l'amour},
qu'on attribue à Pascal, c'est la beauté du transport impétueux de
l'amour qui est magnifiée, et la passion en général est caractérisée
par la précipitation de toutes les pensées dans une seule direction.
Ici s'annonce cette passion de la foi qui s'emparera de lui après sa
conversion définitive. Sa nature garda son caractère primitif à
travers toutes les phases de son évolution religieuse, et bien qu'il
lui fallût passer par une double conversion avant d'en arriver au
terme de son développement, Pascal n'a jamais été un homme sans
religion. Pendant toute sa jeunesse, son âme fut dominée par un
sentiment de gravité religieuse ; mais, à côté, une passion intellectuelle,
d'une activité précoce, l'engageait dans des recherches
qui ont fait époque dans l'histoire des mathématiques et de la
physique ; d'ailleurs cela ne l'empêcha pas quelque temps d'aller
dans le monde et d'y chercher les agréments spirituels et littéraires
du commerce des honnêtes gens. C'est à cette période mondaine
qu'il faut rapporter l'aventure qui a probablement inspiré le \emph{Discours
sur les passions de l'amour}. D'après M. Strowski (\emph{Pascal et
son temps}, II, p. 278) l'objet de sa passion paraît avoir été une
belle dame, de haute naissance, d'humeur aventureuse et passionnée,
mais qui ne lui a point donné matière à repentir. Il n'en devait pas
moins, dans la suite, s'imputer gravement à péché d'avoir eu l'âme
remplie d'autre chose que de Dieu. Car, aux yeux de Pascal, il y
avait là une mortelle antinomie, un antagonisme équivalent au
conflit de la vie et de la mort. Il était de ces natures à qui il faut
une rupture totale avec le passé pour arriver à leur position définitive.
Son « humeur bouillante » le voulait. La catastrophe vint
dans cette nuit du 23 novembre 1654 où, dans une extase, il com%
% At the foot of p. 225, below the text and outside it, the printer's
% signature line:  Rev. Méta. --- T. XXX (n° 2, 1923).        15
% (small caps, the "15" at the outer edge). Not part of the article; not set.
% --- p. 226 ---
prit que le Dieu vivant demande tout et qu'un renoncement plénier
devait s'ensuivre. Sa vie jusqu'à ce jour lui sembla dépensée dans
le vide, maintenant il exultait d'avoir retrouvé la source de vie.

Ayant passé par de tels états psychologiques, Pascal ne pouvait
que s'indigner devant les tentatives que faisaient les Pères de la
Société de Jésus pour adapter aux besoins des gens du monde les
hautes exigences de l'idéal chrétien. Impossible de « joindre Dieu
au monde » (fr. 935 Br.). Le monde ne repousse pas la religion,
mais il la veut tout indulgente et douce (fr. 956 Br.). Dans ses
\emph{Lettres écrites à un provincial}, comme dans les toutes dernières
des \emph{Pensées}, il parle, en termes toujours plus rudes, des voies de
corruption où s'est enfoncée l'Église. « Si elle absout ou si elle lie
sans Dieu, ce n'est plus l'Église » (fr. 870 Br.).

Au cours de la polémique ainsi engagée, dans le temps même où
il publiait les \emph{Provinciales}, il arriva qu'une jeune parente de Pascal
guérit d'une maladie des yeux pourtant déclarée incurable par les
médecins : on lui avait appliqué, sur l'œil malade, une épine de la
couronne du Sauveur, qui faisait partie des reliques de l'église de
Port-Royal. Ce miracle le confirma dans sa conviction d'avoir opté
pour le vrai. L'heure des miracles n'était donc point passée, comme
le soutenaient les adversaires. Et si les autorités de l'Église renient
la vérité, il faut que Dieu lui-même intervienne (fr. 832, 839 Br.).
Aussi Pascal en appelle-t-il avec confiance, de Rome, qui le condamne,
au tribunal de Jésus-Christ (fr. 920 Br.).

Dans sa vie privée, Pascal réalisa le renoncement complet à lui-même.
Sans entrer au couvent, il vécut ses dernières années dans
l'ascétisme le plus rigoureux, tout adonné aux œuvres de charité
et faisant défense inquiète aux siens de s'attacher à sa personne et
de frustrer ainsi Dieu de cet abandon total de l'être qui lui est dû.
La longue maladie qui le rongeait, il la prenait comme un soulagement,
bien qu'elle l'empêchât d'achever cette apologie du christianisme
dont l'ébauche tient dans ses \emph{Pensées}. Pour le chrétien
l'état normal est un état morbide de l'âme où il lui est moins
malaisé de se maintenir quand il souffre aussi dans son corps.

\bigskip
% ^ The journal sets a blank line here --- visibly more than a paragraph
%   break, less than a section break --- dividing the Pascal portrait from
%   the Kierkegaard portrait. No ornament, no rule, no numeral.

C'est un autre portrait moral qui se dégage pour nous de Kierkegaard.
Son évolution a plus de continuité, tout en offrant des
tournants décisifs. Il n'appartient pas à la catégorie des « deux
fois nés ». L'extase ni le miracle ne jouent de rôle dans sa vie.

% --- p. 227 ---
Comme Pascal, il a été élevé dans un christianisme plein de
sérieux, mais durant sa jeunesse il s'occupa de poésie et de philosophie.
C'était le temps du Romantisme et l'idée d'une harmonie
embrassant l'art, la science, la religion, dominait les esprits.
L'idéologie romantique venait d'être mise par Hégel en système et
% ^ "Hégel" with an acute accent, as printed. Ten lines further down the
%   SAME page the name is set "Hegel", unaccented. Both readings verified;
%   the inconsistency is the journal's and is left standing.
Kierkegaard s'y intéressa quelque temps. Dès lors que les antinomies
finissent toutes par se résoudre en harmonie, on pouvait
bien d'abord laisser libre essor à sa rêverie sentimentale et à son
imagination ; ne font-elles point partie, en effet, de ce monde riche
et changeant, dont l'énigme, en fin de compte, se laissait résoudre
par les vérités essentielles du christianisme ? Mais, de bonne heure,
Kierkegaard s'est aperçu qu'à ce jeu on n'échappe pas à une dualité
irréductible : c'est dans ce sens qu'on le voit relever alors, dans
son \emph{Journal}, comme une devise, le vers de Gœthe parlant de
Gretchen : « Halb Kinderspiel, halb Gott im Herzen ! » (Les jeux
% ^ the German is set in ROMAN inside the guillemets, not in italic.
d'enfant et Dieu voisinent dans son cœur). Il a même eu alors une
période où il se laissa entraîner dans une vie de déréglements qui
% ^ "déréglements" with an acute, as printed (mod. "dérèglements").
devait, après coup, nourrir de longs remords cette mélancolie
inhérente à sa nature et par moments si puissante qu'il se sentait
aux bords de la folie.

Il se donna désormais pour mission de fondre en une harmonie
personnelle les poussées contradictoires de son être --- insouciance
et gravité, frivolité et mélancolie. L'expérience personnelle lui fit
voir alors qu'il était autrement malaisé en pratique d'atteindre à
cette « unité supérieure » que réalise \emph{de plano} le système de Hegel.
Ce travail intérieur fit de lui un solitaire. Il se sentait isolé et mal
compris. Sa pente mélancolique lui interdisait, pensait-il, d'entrer
dans des liens impliquant un abandon sincère et joyeux. Il rompit
ses fiançailles avec une jeune fille d'humeur gaie et enjouée et il se
jugeait incapable d'exercer une fonction quelconque. Le contraste
était trop criant, à ses yeux, entre l'aisance légère avec laquelle les
autres prenaient l'existence et la lutte où il se débattait contre les
forces assombrissantes de sa vie. Il se disait qu'au moyen-âge son
entrée au couvent eût été tout indiquée. Cette obscure et patiente
élaboration de sa vie personnelle, il l'a décrite sous cette image
poétique : « Je suis aux écoutes de mes musiques intérieures, des
appels joyeux de leur chant et de leurs basses notes graves d'orgue.
Et ce n'est pas petite tâche de les coordonner quand on n'est pas un
organiste, mais un homme qui se borne, à défaut d'exigences plus
% --- p. 228 ---
grosses envers la vie, au simple désir de se vouloir connaître ». ---
Qu'on est donc loin du temps où « Dieu » et les « jeux d'enfant »
voisinaient en paix dans son cœur. Leur antagonisme constitue un
problème dont la solution demande l'énergie tout entière de la vie.
Et, en effet, l'effort pour concilier les inconciliables est caractéristique
de toute la vie de Kierkegaard. Il ne va pas d'un extrême à
l'autre : les contraires sont là, en présence, depuis toujours, seulement
leur opposition se fait plus tranchante au fur et à mesure
qu'avance le développement de son individualité, et l'effort pour les
vaincre s'en augmente à proportion. C'est à l'échelle de sa propre
expérience que Kierkegaard a pu mesurer l'intensité d'effort de la
vie personnelle.

Un moyen s'offrait à lui pour tenir la mélancolie en échec et soulager
la tension intérieure : il n'avait qu'à se laisser aller à son
penchant à la production littéraire. Ce lui était une délivrance de
s'adonner à la production littéraire.
% ^ The phrase "à la production littéraire" is set TWICE in two consecutive
%   sentences. Almost certainly a compositor's dittography, but it is what
%   the page carries and it is left exactly as printed.
Comme la Princesse des \emph{Mille
et une nuits} qui contait pour sauver sa vie, il sauva la sienne,
dit-il, en écrivant. A la première période de l'écrivain (1843-1846)
remontent quelques-uns de ses écrits les plus connus : \emph{Ou l'un ou
l'autre} et \emph{Étapes de la route humaine}. Ils roulent sur l'objet
même de ses luttes intimes : l'élaboration intérieure de la personnalité,
la conquête d'un noyau solide autour duquel graviteront les
autres éléments de l'âme. Par là ses premières œuvres purent servir
en outre à stimuler la génération énervée du Romantisme et de
la Philosophie spéculative que les épigones entraînaient à émousser
les antagonismes de l'homme et à rabaisser de ce fait sa vie spirituelle.
Kierkegaard n'a pu prétendre qu'il ait eu en vue d'exercer
une telle influence à l'origine de sa carrière. Honnêtement il se rendait
compte qu'il y avait eu, en lui, une tension qui appelait un
dérivatif. Au contraire de Pascal qui se proposait délibérément,
dans ses \emph{Provinciales}, de combattre un affaissement de la morale,
Kierkegaard en vint peu à peu, et sans l'avoir vraiment voulu, à
travailler dans le même sens. Il n'était d'ailleurs pas encore à ce
moment aussi avant que Pascal dans son développement religieux.

Une nouvelle période dans la carrière de Kierkegaard (1847-1855)
s'ouvrit justement parce que sa vie intérieure prit décidément un
caractère de plus en plus religieux. Il écrit dans son \emph{Journal}
(1847) : « J'éprouve maintenant le besoin d'une compréhension de
plus en plus profonde de moi-même en me rapprochant toujours
% --- p. 229 ---
plus de Dieu. Il s'ébauche en moi je ne sais quoi qui annonce une
métamorphose... Il faut donc que je me tienne tranquille. » On
remarquera le contraste avec l'extase de Pascal : ici, l'éclosion
presque inconsciente d'un nouveau mode d'existence, devant
laquelle le sujet garde une attitude passive. L'année d'après, il
note : « A présent, je possède la foi dans l'intime acception du
terme ». Le christianisme prenait à ses yeux un caractère de réalité
que jusqu'alors il ne lui connaissait pas. Cette première
période de sa carrière terminée, il crut quelque temps que son
action littéraire aussi avait pris fin ; l'idée lui souriait de se retirer
dans un coin perdu à la campagne. Mais il ne la mit point à exécution :
la force accrue de sa conscience religieuse éveilla une critique
plus sévère de la chrétienté existante. Dans sa première
période littéraire, Kierkegaard avait flétri le relâchement où était
tombée la conception personnelle de la vie ; le chrétien, maintenant,
découvrait dans l'Église une facilité de transaction, un esprit d'accommodement
du christianisme aux goûts des mondains, qui reléguaient
à l'arrière-plan l'idéal des premiers chrétiens. C'eût été,
d'après lui, supprimer virtuellement le christianisme que d'en
faire une religion toute de douceur et de consolation. « Et nous
voyons vivre, dans la chrétienté actuelle, une génération gâtée,
fière, et lâche pourtant, arrogante mais sans ressort, qui se laisse
administrer de temps en temps ces bons principes consolatoires,
sans même savoir au juste si elle en usera quand la vie lui sourit,
et qui s'en scandalise aux heures de détresse quand il appert qu'au
fond leur indulgence se dérobe. » Nous empruntons cette citation
à l'\emph{Exercice dans le Christianisme}, qui caractérise, avec la \emph{Maladie
à la mort}, cette deuxième période de la production de Kierkegaard.

L'Église officielle ne releva pas la constatation d'un conflit éclatant
entre l'obligation âprement imposée par le christianisme primitif
de ne poursuivre que notre unique nécessité et, d'autre part,
l'idylle organisée par la chrétienté moderne sous le couvert du
dogme de la Rédemption. C'est alors que Kierkegaard entama sa
guerre passionnée contre l'Église existante, la plus violente de
toutes les luttes qu'ait connues notre histoire intellectuelle danoise.
Le dernier mot de Kierkegaard y fut : « Le christianisme du Nouveau
Testament n'existe pas ». En plein combat la maladie l'arrêta
et la mort.
% p. 229 breaks off short --- roughly four lines of white below the last
% line --- so section I ends here and section II opens on p. 230.

% Mechanical OCR skeleton for pp. 230--238, built by build_frag.py.
% Words from tesseract -l fra; structure from TSV geometry.
% STILL TO DO BY EYE, against the page images:
%   italics (OCR cannot see them), section heads, footnote placement,
%   guillemet spacing, printer's defects.

% --- p. 230 ---

% Section II opens here, set exactly as section I on p. 225:
% the numeral and dash in full capitals, then caps + small caps.
\section*{II. --- \textsc{Prédispositions intellectuelles.}}

Pascal poursuivit ses études de mathématiques et de physique et
s'y est acquis un beau nom. C'est à ces études qu'il dut sa conception
de la connaissance scientifique. Un système comme celui de
Descartes, tendant à réduire la connaissance de la nature à certaines propositions fondamentales et à préparer ainsi des bases
rationnelles au travail de la raison humaine --- un tel système lui
semblait condamné par avance. On lui a appliqué justement le
nom de positiviste expérimentaliste. Dès sa période « mondaine »
il avait compris qu'étant données la multiplicité, la complexité et la
variété des phénomènes psychiques, ce n'est pas à « l'esprit de
géométrie », c'est-à-dire au raisonnement procédant par déductions
purement mathématiques qu'il faut demander la compréhension de
ce qui se passe dans l'âme humaine. En psychologie, il importe
avant tout de rapporter à certains types les phénomènes multiples,
de s'identifier par la pensée à chacun de ces types et de chercher à
en connaître les mobiles et les préoccupations. Ici le raisonnement
mathématique se trouverait en défaut ; seul « l'esprit de finesse »,
sens subtil des types et des nuances, nous permettra de démêler
les éléments typiques et de les retrouver dans leurs combinaisons
variées. M\textsuperscript{me} Perier admirait précisément chez son frère son aptitude
à se rendre compte des données mentales de ses interlocuteurs.
Il connaissait, pense-t-elle, « tous les ressorts du cœur
humain ». Cette faculté avait été développée en lui par la fréquentation
des gens du monde ; elle s'est montrée utile dans la suite,
quand il s'est agi pour lui de gagner des âmes pour la cause de la
vraie Religion. Les preuves fournies par la spéculation sont peu
efficaces. Concluantes ou non, leur action se borne aux moments
où nous subissons l'influence des mouvements intellectuels ;
ensuite elles sont facilement oubliées. La bonne tactique, c'est de
prendre les hommes au point où ils en sont restés et de leur
montrer à aller de l'avant par leurs propres ressources : rien ne
convainc mieux les hommes que les raisons qu'ils tirent de leur
propre fonds. Il faut les aider à voir le mélange curieux de grandeur
et de misère qu'est leur nature. La pensée pourra bien les enlever
dans ses hautes régions, mais la concupiscence et les passions
les ramèneront toujours vers la fange d'en bas. L'homme doit être
amené à comprendre quel « monstre incompréhensible » il fait, et
% low-confidence on p.230: II., INTELLECTUELLES.
% low-confidence on p.230: “,  taines
% low-confidence on p.230: qu'étant
% low-confidence on p.230: l'esprit
% low-confidence on p.230: c'est-à-dire
% low-confidence on p.230: de chercher
% low-confidence on p.230: préoccupations.
% low-confidence on p.230: l'esprit
% low-confidence on p.230: démêler
% low-confidence on p.230: combinaisons
% low-confidence on p.230: M”°
% low-confidence on p.230: mentales, interlocu-
% low-confidence on p.230: I]
% low-confidence on p.230: Cette
% low-confidence on p.230: l'influence
% low-confidence on p.230: “ensuite
% low-confidence on p.230: convainc, qu'ils
% low-confidence on p.230: qu'est
% low-confidence on p.230: et, les
% low-confidence on p.230: d'en
% --- p. 231 ---
seule la vraie Religion lui fournira la clef de sa nature avec les voies
et moyens de résoudre les divisions intimes qui le déchirent. Pascal
tient surtout à ce que l'homme arrive à se sentir seul dans un
monde infini d'où n'émane aucune voix de compassion ou d'encouragement.
Le trouble, le désarroi, l'angoisse qu'éprouvait Pascal
en face des grandes énigmes lui fit prendre la seule voie qui s'ouvrait,
l'issue du christianisme. Sans doute, le christianisme contient
des mystères à côté des révélations qu'il apporte aux hommes.
Mais si Dieu ne nous avait rien laissé ignorer du mystère de notre
vie dans l'univers, n'eût-il pas du même coup supprimé dans l'âme
cet effort, ce combat, ce courage d'oser et de jouer son va-tout, qui
lui permettent seuls de mériter l'éternité ?

Pascal ramène à deux types principaux les hommes qui tentent
de résoudre, par les lumières naturelles, les grandes énigmes de
l'existence. Les uns s'attachent à la raison pensante et au devoir
qui sont proprement la grandeur et le point solide dans l'homme,
et le reste ils le tiennent pour négligeable. Leur représentant classique,
c'est Épictète le Stoïcien. Les esprits de l'autre type s'accommodent
des fluctuations de la vie, ils pratiquent le \emph{carpe diem} et
s'en tiennent d'ailleurs aux opinions reçues, se rendant bien compte
que tout est sujet à vicissitude. Ils se reposent avec délice dans la
certitude de n'être sûrs de rien. Montaigne est le parangon des insouciants
conscients.

L'étude de ces deux types fut pour Pascal une préparation au
corps à corps où il se proposait, dans ses \emph{Pensées}, d'empoigner et
de secouer l'incrédule, afin de le mettre dans l'état critique où il
n'apercevrait qu'une seule et unique porte de salut --- celle qu'avait
trouvée Pascal lui-même. Ce lui fut une déception de voir combien
était peu appréciée l'étude de la nature humaine, la recherche de
ses types et de leur utilité pratique. « Quand j'ai commencé l'étude
de l'homme, j'ai vu que les sciences abstraites ne sont pas propres
à l'homme, et que je m'égarais plus de ma condition en y pénétrant
que les autres en les ignorant. » Il avait espéré, alors, trouver des
compagnons plus nombreux en l'étude de l'homme, mais ici, contre
son attente, il constata plus de détachement encore que dans le cas
de la géométrie (fr. 144 Br.). Lui-même se jeta à corps perdu dans
la nouvelle étude se désintéressant de plus en plus de ses premiers
travaux scientifiques. Dans une lettre au célèbre mathématicien
Fermat, il écrit vers la fin de sa vie (10 août 1660) que, tout en con\-%
% low-confidence on p.231: eb
% low-confidence on p.231: l'angoisse
% low-confidence on p.231: s'ou-
% low-confidence on p.231: l'univers,
% low-confidence on p.231: d'oser, qui, _----
% low-confidence on p.231: l'éternité
% low-confidence on p.231: l'homme,
% low-confidence on p.231: s'accom-
% low-confidence on p.231: opinions, recues,
% low-confidence on p.231: est, le
% low-confidence on p.231: ‘
% low-confidence on p.231: L'étude
% low-confidence on p.231: d'empoigner
% low-confidence on p.231: n'apercevrait, ;
% low-confidence on p.231: l'étude
% low-confidence on p.231: l'homme,
% low-confidence on p.231: Il
% low-confidence on p.231: l'étude
% low-confidence on p.231: dans
% low-confidence on p.231: :
% --- p. 232 ---
sidérant toujours la géométrie comme le domaine le plus élevé où
puissent s'exercer les facultés de l'esprit, il la regarde désormais
comme tout à fait inutile, surtout depuis qu'il a entrepris des études
tellement différentes de ses travaux antérieurs que ces derniers se
sont presque évanouis de sa mémoire.

Il reste convaincu que la grandeur et les titres de noblesse de
l'homme résident dans la pensée et que c'est à elle qu'il faut nous attacher.
Comparés aux grandes masses qui nous entourent nous
ne tenons que peu d'espace et notre puissance est négligeable auprès
de la leur. Nous sommes des roseaux, mais des roseaux pensants,
et les masses extérieures ont beau nous écraser, elles n'atteindront
jamais la dignité que confère la pensée (fr. 346-348). Mais cette
pensée tant exaltée par Pascal se réduisait de plus en plus, pour
lui, à n'être qu'une servante de l'explication théologique. De même
que, pendant sa période « mondaine », il avait décrit la passion de
l'amour comme consistant dans une concentration, une précipita\-%
tion de toutes les pensées vers un seul but, de même, maintenant
que la religion était devenue sa passion dominante, il n'admettait
que des pensées orientées dans cette seule direction.

Dans le cas de Kierkegaard, les données intellectuelles sont tout
autres. Grandi parmi les épigones du Romantisme, il vit, pendant
quelques années, dans la métaphysique hégélienne, l'idéale philosophique,
tandis qu'aux yeux de Pascal la pensée idéale se trouvait
représentée par la géométrie. Dans un système comme celui de
Hegel, où une nécessité intérieure enchaîne toutes les parties les
unes aux autres, il n'est rien qui n'ait sa place logiquement mar\-%
quée. La Poésie, la Religion, la Science n'étaient donc pour Kierkegaard
que les formes diverses ou successives d'une même et
divine vérité. Tel fut le rêve de sa jeunesse, dont il finit enfin par
s'affranchir --- sur les pas de son maître, le professeur Paul Moller,%
% As printed: "Paul Moller", for Poul Møller. Høffding's own French
% form of his teacher's name, without the ø. Not corrected.

le plus danois de tous les écrivains danois. Le mérite de Kierkegaard,
dans l'ordre philosophique, est d'avoir été un des premiers à
critiquer la philosophie spéculative du point de vue de la théorie
de la connaissance (\emph{Postscriptum définitif non scientifique}, 1846)%
% printed mark: superscript 1 followed by a baseline ')'
\footnote{\dk{\emph{Alsluttende uvidenskabelig Efterskrift.}} Copenhague, 1846.}%
% PRINTER'S ERROR, transcribed as printed: the journal sets
%   "Alsluttende" for Kierkegaard's "Afsluttende". Verified at 600 dpi.
.
Il montre que même s'il existait une connaissance absolue, chacun
de nous serait obligé pourtant de commencer, à un point déterminé,
% low-confidence on p.232: ,
% low-confidence on p.232: 1
% low-confidence on p.232: Il
% low-confidence on p.232: nous, -
% low-confidence on p.232: d'espace
% low-confidence on p.232: pour
% low-confidence on p.232: n'être
% low-confidence on p.232: ,
% low-confidence on p.232: {
% low-confidence on p.232: {
% low-confidence on p.232: {
% low-confidence on p.232: -
% low-confidence on p.232: scientifique,, 1846)*,
% low-confidence on p.232: 1,, Alsluttende, uvidenskabelig, Efterskrift.
% --- p. 233 ---
pris dans le domaine de l'expérience, le raisonnement qui devait
nous élever à ce stade supérieur, et pourquoi, alors, à tel point
plutôt qu'à tel ou tel autre ? Et ce point de départ, choisi dans le
monde réel, pourra-t-on le déduire à son tour de la connaissance
absolue qui doit en résulter ? Et ne subsistera-t-il pas toujours
d'autres données que ce point de départ expérimental à la base
même de la connaissance absolue, et quelle certitude aura-t-on
d'être remonté jusqu'aux plus éloignées ? Kierkegaard a été amené
à voir que de tous les problèmes philosophiques celui de la connaissance est le plus central. Il insiste sur le fait que nous vivons
dans le temps. Nous vivons en avant, mais nous comprenons en
arrière. De là, l'impossibilité d'un aboutissement terminal, et Kierkegaard
se déclare d'accord avec Lessing disant que pour l'homme
il n'y a rien au delà de la recherche de la vérité.

Comme Pascal, Kierkegaard est d'avis que la vie personnelle a
besoin de pensées, mais de pensées appliquées à ce qui étaye notre
existence. La pensée de Dieu n'est qu'un expédient désespéré
et n'offre pas d'objet à l'intelligence pure. Kierkegaard établit une
distinction entre les connaissances « essentielles » et les connaissances
« non essentielles », en classant parmi les connaissances non
essentielles : sciences naturelles, philologie, histoire et philosophie
pure --- toutes pensées sans emploi pour nous quand nous prenons
position dans le problème capital de la vie et de la mort. Impossible
de concevoir la vie d'une façon purement intellectuelle, à moins de
faire de soi-même une abstraction. Pour apprécier les diverses attitudes
que prennent les hommes vis-à-vis de l'existence, Kierkegaard
se sert d'un procédé qui rappelle celui de Pascal : il les groupe
dans une série de types qu'il entreprend de bien caractériser. Ne
retenons ici que ceux dont il a tracé le portrait dans ses écrits les
plus populaires (\emph{Ou l'un ou l'autre} et \emph{Étapes de la route humaine}).
L'« esthéticien » veut jouir par-dessus tout du moment qui passe.
Son art de vivre, c'est de garder assez d'élasticité d'âme pour saisir
la nouveauté successive des instants. Il importe de couper court au
moment opportun, afin de ne pas être retenu dans d'immobiles
liens. La méthode des cultures alternées a du bon. L'esthéticien est
comme une tangente à la circonférence (plus exactement : aux circonférences)
de la vie. À un seul point il y a, momentanément, un
contact, qui est aussitôt suivi par un contact à un autre point. Donc ;
des amours qui changent et pas de mariage ; des camaraderies et
% low-confidence on p.233: jusqu'aux
% low-confidence on p.233: :
% low-confidence on p.233: d'accord
% low-confidence on p.233: ,
% low-confidence on p.233: d'avis, personnelle
% low-confidence on p.233: ‘de, ;
% low-confidence on p.233: d'objet
% low-confidence on p.233: philosophie
% low-confidence on p.233: moins, de
% low-confidence on p.233: :
% low-confidence on p.233: :
% low-confidence on p.233: l'un, ou
% low-confidence on p.233: court
% low-confidence on p.233: “, cir-, -
% low-confidence on p.233: il!
% --- p. 234 ---
pas d'amitiés ; des occupations improvisées et pas de devoirs professionnels ;
des jouissances intellectuelles et point de science ! Si
la vie « esthétique » est une tangente, la vie « éthique », elle, décrit
une circonférence autour d'un centre et se trouve déterminée
par rapport à ce centre. C'est le type de la fidélité. Sa vie est une
vie en commun, une vie réglée, une vie de travail incessant. Au-dessus
de ces deux types s'élève celui de la vie religieuse dont
l'effort et l'orientation se rapportent à un but situé au delà de tous
moments et de toute durée, incommensurable à tous les autres buts
humains et même, à son point culminant, dans le christianisme,
en contradiction manifeste avec tout ce qui, d'un point de vue purement
humain, pourrait être pris comme but et comme directive
de l'existence. Plus se différencie le caractère de la vie religieuse,
moins elle permettra à l'homme de se développer dans une
ambiance naturelle : il sera comme un poisson sur terre. Finalement,
l'Éternité ne restera pas (comme chez Platon et chez Spinoza)
le fond immense, toujours présent sur lequel se déroule la vie
mesurée dans le temps : elle fera irruption dans la durée, le divin
surgissant sous forme humaine à un point déterminé du temps et
de l'espace. Alors, cette révélation de l'attitude des hommes à son
égard sera le facteur décisif pour apprécier la vie et son contenu.

C'est une philosophie comparée de la vie que nous donne ici
Kierkegaard. Elle fait penser aux caractères de Montaigne et d'Epictète,%
% As printed: "Epictète" without the acute here, though the same name
% is set "Épictète" on p. 231. Both readings recorded, not normalised.

décrits par Pascal, et à celui du chrétien qu'il leur oppose à
tous deux. Et de même que Pascal voyait dans la maladie l'état
naturel aux chrétiens, de même Kierkegaard, de son côté, aboutit
à ce résultat que le type de vie, le plus élevé, celui des chrétiens,
comporte tension et souffrances dues à l'antagonisme déchirant des
contradictoires entre lesquels la vie doit être vécue. Et la supériorité
de ce dernier type sur les autres vient justement des plus
grandes épreuves qu'il impose à cette puissance de synthèse qu'exige
toute vie personnelle.

D'après Kierkegaard, c'est par bonds qu'on passe, comme de
stade en stade, d'une vie à l'autre, et le bond principal est celui qui
porte une âme plongée dans le désespoir à la vie en Christ. Mais,
même au point le plus décisif, ce bond ne prend jamais, chez
Kierkegaard, le caractère de l'extase. Et, dans une certaine mesure,
la continuité subsiste, puisque notre énergie synthétique reste
l'échelle commune à tous les types de vie personnelle. D'ailleurs
% low-confidence on p.234: occupations, pro-
% low-confidence on p.234: ecrit
% low-confidence on p.234: ;, rapport
% low-confidence on p.234: l'orientation
% low-confidence on p.234: religieuse,
% low-confidence on p.234: sur, |
% low-confidence on p.234: ici
% low-confidence on p.234: d'Epic-
% low-confidence on p.234: celui
% low-confidence on p.234: °
% low-confidence on p.234: l'antagonisme
% low-confidence on p.234: {
% low-confidence on p.234: #, qu'il
% low-confidence on p.234: :
% low-confidence on p.234: qu'on, comme
% --- p. 235 ---
Pascal emploie au fond la même mesure ; il la tire du rapport entre
le fonds plus ou moins riche qu'offre une existence et l'énergie de
l'homme à concentrer et à ordonner ses richesses. Aussi bien est-ce
l'unique mesure qu'on puisse employer vis-à-vis des différences,
des individus et des époques, dans leur manière de comprendre
la vie ; il s'impose même à ceux qui ne placent pas le but final de
l'existence où l'ont mis Pascal et Kierkegaard. Une objection nous
vient dès maintenant contre cet emploi qu'ils font de leur principe
de mesure. Ils considèrent de façon exclusive la résistance intérieure,
subjective qu'il faut vaincre, par conséquent, le travail intérieur par lequel s'obtient l'unité et la concentration de l'âme. Mais
la question se pose de savoir quelle est l'importance de ce travail
intérieur pour les groupes humains où vit l'individu, pour la
société dont il fait partie, peut-être même pour l'humanité entière.
C'est là un point de vue qu'on ne ferait jamais adopter à Pascal ni
à Kierkegaard. D'après eux, il y a tout un abîme entre le type de
vie qu'ils regardent comme le plus élevé et les autres vies consacrées à l'œuvre de la civilisation humaine.

% Section III opens here.
\section*{III. --- \textsc{Le problème chrétien.}}

La vie et l'œuvre de Pascal s'accomplirent à l'intérieur du catholicisme : Kierkegaard a évolué dans un milieu protestant. De là,
entre eux, des différences multiples se manifestant dans leur formation
religieuse aussi bien que dans les positions où ils se sont
définitivement établis. Mais, au point de vue de l'histoire de la religion,
il est très intéressant de voir qu'arrivés au terme de leur
effort pour découvrir le point central de leur religion, ils s'arrêtent
devant un seul et même problème, qu'on pourrait formuler ainsi :
Est-il possible de rester dans l'esprit et la pratique du christianisme
primitif, % A single damaged/over-inked sort stands immediately before
% "s'assimilant" here; it is not legible as any letter at 600 dpi.
% Sense wants "tout en se l'assimilant". Left as the visible text.
tout en s'assimilant et en servant une culture, morale
et matérielle, qui non seulement ne sort pas du christianisme, mais
même l'a précédé et a continué de se développer en dehors de lui
d'après ses lois particulières ? En effet, le problème se pose pour
les catholiques comme pour les protestants. Avant d'aller plus loin,
considérons ce qui l'a fait naître.

Les premières générations chrétiennes vécurent dans l'attente
d'un retour prochain du Christ, qui devait clore les destinées du
monde par un jugement universel. Au début, cette attente constituait
la croyance des chrétiens ; elle réglait leur vie et leur concep\-%
% low-confidence on p.235: [
% low-confidence on p.235: et
% low-confidence on p.235: l'unique
% low-confidence on p.235: s'impose
% low-confidence on p.235: l'existence
% low-confidence on p.235: qu'ils
% low-confidence on p.235: vaincre,, {
% low-confidence on p.235: l'individu,
% low-confidence on p.235: l'humanité
% low-confidence on p.235: abime
% low-confidence on p.235: {
% low-confidence on p.235:  eérées
% low-confidence on p.235: III.
% low-confidence on p.235: És'assimilant
% low-confidence on p.235: d'aller
% low-confidence on p.235: du
% --- p. 236 ---
tion de la vie. Dans la perspective d'une catastrophe imminente
d'où naîtraient un nouveau ciel, une terre nouvelle, les devoirs
purement humains ou sociaux perdaient leur urgence. C'est pour\-%
quoi il n'est pas question, dans le Nouveau Testament, d'art, de
science, de vie publique ; la vie de famille elle-même n'était admise
que sous certaines réserves. La morale du Nouveau Testament devait
être ce que M. Albert Schweitzer a bien nommé une « éthique intérimaire
». Tous les efforts, tous les projets convergeaient sur un
unique but sis dans un avenir voisin. Pour cette chrétienté primi\-%
tive, l'avènement du règne de Dieu, auquel il est fait allusion dans
la prière dominicale, prenait un sens très précis : on n'entendait
point par là demander vaguement l'obtention de biens spirituels.
Chez les fervents, l'attente se doublait d'un enthousiasme haussé
souvent jusqu'à de l'extase et qui rendait faciles tous les sacrifices.
Les choses de la vie journalière des hommes se rapetissaient devant
la grandeur de celles qu'on attendait.

À l'encontre d'une croyance comme celle des chrétiens primitifs,
le catholicisme et le protestantisme apparaissent comme des organisations
à longue échéance basées sur la chance d'une évolution
prolongée de la vie religieuse et, en général, de la vie humaine sur
terre. De moins en moins, on s'est préparé à un départ brusqué.
Tout en s'efforçant de ne pas sortir de l'attente des premiers siècles,
on en a placé cependant la réalisation dans un avenir toujours plus
reculé, la sujétion où elle avait tenu les âmes allait en diminuant.
Peu à peu on s'appropria, --- après une résistance convenable, qui
fut souvent assez âpre, --- non seulement la culture transmise par
l'antiquité classique, mais aussi la culture nouvelle qui se développait
indépendamment du christianisme, fait qui n'empêche point
les deux Églises de s'affirmer en continuité avec la croyance des
premières générations chrétiennes et de maintenir, chacune à sa
manière, la prétention d'avoir gardé intacte la vieille tradition à
travers les temps nouveaux.

La solution la plus grandiose du problème est celle du catholicisme.
Elle se fonde sur la distinction de degrés dans la perfection.
Religieux et religieuses renoncent à la vie dans le siècle afin de
pouvoir se consacrer entièrement à l'au-delà, à la seule chose nécessaire.
C'est par eux que subsiste la filiation du christianisme originel.
Aux laïques incombe la tâche de continuer la vie humaine sous
% low-confidence on p.236: ;, imminente
% low-confidence on p.236: purement
% low-confidence on p.236: |
% low-confidence on p.236: nommé
% low-confidence on p.236: :
% low-confidence on p.236: -
% low-confidence on p.236: doublait
% low-confidence on p.236: |
% low-confidence on p.236: [
% low-confidence on p.236: orga-
% low-confidence on p.236: s'efforçant, l'attente
% low-confidence on p.236: |
% low-confidence on p.236: s'appropria,
% low-confidence on p.236: transmise, [
% low-confidence on p.236: n'empêche
% low-confidence on p.236: s'affirmer
% low-confidence on p.236: |, de., maintenir,
% low-confidence on p.236: ,
% low-confidence on p.236: nouveaux.
% low-confidence on p.236: origi-
% --- p. 237 ---
l'influence des directeurs de conscience par lesquels l'Église exerce
son action de contrôle sur les âmes. Saint Augustin (\emph{De civitate
Dei}, XX, 9) voyait déjà, dans les prélats de l'Église, une réalisation
partielle de la prophétie de l'Apocalypse, parlant de ceux qui
occupent le tribunal dans le royaume de mille ans.

Pascal s'est fort bien rendu compte qu'aux yeux de l'Église, l'éloignement
du christianisme originel pour toutes les formes de culture
ne devait jamais cesser de prévaloir dans une certaine pratique.
Sans entrer au couvent, sa vie des dernières années n'en fut pas
moins d'un moine. À ses yeux, les temps qui nous séparent du
retour du Christ ne doivent être qu'attente ardente et que souf\-%
france. La scène de Gethsémani représente bien, selon lui, la situation
permanente du vrai chrétien. Nous devons vivre en contemporains
actifs des souffrances du Sauveur. Une seule fois, à Gethsémani,
Jésus, en détresse, a demandé du secours, du réconfort aux hommes
et ses apôtres lui firent défaut. Mais le devoir auquel ils manquèrent
alors incombe aux chrétiens des âges successifs. « Jésus
sera en agonie jusqu'à la fin du monde ; il ne faut pas dormir pendant
ce temps-là » (fr. 583 Br.). Donc, ce n'est pas vers le Christ
vainqueur qu'il faut tourner les regards, comme le fait le plus souvent
l'Église, c'est vers Jésus souffrant. Le jardin des Oliviers est la
contre-partie du jardin de l'Éden. L'homme déchu doit faire de
Gethsémani le séjour permanent de ses pensées ; c'est à ce prix
seulement qu'il garde l'espoir de voir s'ouvrir devant lui les portes
de l'Éden.

La direction des âmes telle que la pratiquaient les pères jésuites
apparaissait à Pascal comme un accommodement criminel de la vie
selon le Christ à la vie selon le monde, et, dans les \emph{Lettres à un pro\-%
vincial}, il a élevé contre les abus des jésuites une vive protestation.
Sauf des erreurs de détail, il a raison au point de vue rigoureusement
chrétien. Du Nouveau Testament à la casuistique des jésuites l'éthique
ne reste pas identique, et même d'un point de vue purement humain,
on pourrait émettre des doutes touchant les résultats de cette casuistique.
À côté de sa campagne contre la religion de douceur qui s'était,
selon lui, substituée à la méditation des souffrances de Gethsémani,
Pascal a consacré un mémoire à la \emph{Comparaison des chrétiens des
premiers temps avec ceux d'aujourd'hui}.

S'il y a tant d'écart entre les deux catégories de chrétiens, c'est
parce qu'à présent le baptême précède l'instruction des mystères de
% low-confidence on p.237: :, ,
% low-confidence on p.237: ;
% low-confidence on p.237: foyaume
% low-confidence on p.237: .
% low-confidence on p.237: qu'attente
% low-confidence on p.237: ;
% low-confidence on p.237: jusqu'à
% low-confidence on p.237: L'homme
% low-confidence on p.237: s'ouvrir
% low-confidence on p.237: ;
% low-confidence on p.237: Lettres, ,
% low-confidence on p.237: d'un
% low-confidence on p.237: s'était,
% low-confidence on p.237: °
% --- p. 238 ---
la religion, tandis qu'à l'origine, personne n'avait accès au chris\-%
tianisme sans avoir été bien instruit de sa doctrine et sans avoir
renoncé au monde. C'est qu'alors la distinction s'imposait entre le
« siècle » et l'Église, tandis qu'aujourd'hui l'Église a admis le siècle
dans son sein en accordant le baptême aux enfants. L'institution du
baptême des enfants, ajoute Pascal, a été faite à bonne intention :
on craignait de priver du salut les enfants qui mouraient non baptisés,
mais ce qu'on a fait pour le bien des enfants a tourné à « la perte des
adultes ! »

La rupture de Pascal avec l'Église n'a pas été amenée par la divergence
constatée sur ce seul point. Comme nous l'avons vu déjà, il a
eu pour s'insurger contre l'autorité de l'Église des raisons tirées non
seulement de la pratique de l'Église, mais encore de ses dogmes ; il
s'est agi également de l'établissement des faits. D'accord avec ses
amis jansénistes, Pascal affirmait que les propositions incriminées et
condamnées par le pape comme hérétiques n'étaient pas dans l'ouvrage
de Jansénius intitulé \emph{Augustinus}, et il refusait au pape le droit
d'ériger en fait acquis un point expressément vérifiable par tout le monde.
À une occasion antérieure, il avait déjà, dans son \emph{Mémoire sur
le Vide}, repoussé l'ingérence de l'autorité religieuse dans les pro\-%
blèmes des sciences expérimentales. Et dans sa \emph{18\textsuperscript{e} Provinciale} il expose
que si l'opinion de Galilée sur le mouvement de la terre est une
erreur, ce n'est pas un décret du pape qui le démontrera, et qu'au cas
où il est possible de prouver que c'est la terre qui tourne, elle entraîne
dans son mouvement ceux qui le nient comme ceux qui l'admettent.
Pascal n'était pas copernicien, mais il touche ici à une matière fort délicate
qui avait fait condamner au feu Giordano Bruno et rendu parjure
Galilée. Un même souffle rebelle à l'autorité religieuse fait la liaison
secrète entre le Pascal physicien et le Pascal des \emph{Pensées}. Il va jusqu'à
dire que « le pape hait et craint les savants qui ne lui sont pas
soumis par vœu » [comme le sont les religieux et les prêtres] (fr. 873,
Br.). La remarque visait la condamnation par le pape de la doctrine
augustinienne de la grâce efficace et de l'absolue dépendance de la
volonté humaine de par rapport à la volonté de Dieu, après que
Jansénius, en la reprenant, avait accentué l'opposition à la théologie scolastique et jésuite. Le pape lui-même ne serait donc pas orthodoxe,
et son autorité n'est pas plus infaillible en matière de foi que
lorsque les faits sont en cause. C'est que le Pape doit veiller à tant
d'affaires qu'il n'a pas le temps de s'occuper des choses de la foi
% low-confidence on p.238: ‘
% low-confidence on p.238: qu'alors
% low-confidence on p.238: qu'aujourd'hui
% low-confidence on p.238: ,
% low-confidence on p.238: fait, pertedes
% low-confidence on p.238: ;, :
% low-confidence on p.238: l'autorité
% low-confidence on p.238: l'établissement
% low-confidence on p.238: quelles
% low-confidence on p.238: :
% low-confidence on p.238: ilavait
% low-confidence on p.238: l'autorité
% low-confidence on p.238: :
% low-confidence on p.238: qui, le
% low-confidence on p.238: avait
% low-confidence on p.238: l'autorité
% low-confidence on p.238: qu'à
% low-confidence on p.238: ‘
% low-confidence on p.238: {
% low-confidence on p.238: après
% low-confidence on p.238: logiescolastique, et, jésuite.
% low-confidence on p.238: d'affaires
% --- p. 239 ---
(fr. 882). Mais si Rome condamne la vérité, il faut crier d'autant plus
fort ; il faut en appeler au ciel. « Il est meilleur d'obéir à Dieu qu'aux
hommes » (fr. 920). Les amis jansénistes de Pascal jugèrent bon,
par la suite, de supprimer ces dernières lignes dans la première édition
des \emph{Pensées}. Ce lui fut un grand chagrin de sentir que sur ce point
ils ne mettaient pas la vérité au-dessus de tout. La mort seule l'avait
empêché de publier lui-même ses dernières pensées. A son lit de mort
% ^ "A son lit de mort" --- capital A unaccented, as printed.
il reçut les sacrements de la main de son curé ; ses adversaires l'interprétèrent
comme une rétractation, mais il y faut voir seulement
qu'il ne pensait pas avoir rompu avec la vraie Église.

Kierkegaard a grandi au sein de l'Église protestante, dans la vénération
de ses chefs, contemporains et anciens. C'est pendant les
années 1847-1850 qu'il aperçut de plus en plus le désaccord qui
existe entre le christianisme primitif et celui des temps modernes, ---
% ^ The em-dash at the end of this line is genuine printed type (a solid
%   48 px bar at 300 dpi, em-dash width, on the line's own baseline and
%   inside the measure), but the compositor filled the line with about four
%   word-spaces of quads before it: the other word gaps on the line run
%   20--34 px, this one 120 px. Comma + em-dash is this journal's habit
%   (cf. "dans son \emph{Journal}, ---" on p. 224).
et deux de ses plus pénétrants écrits, \emph{La Maladie à la mort} (1849) et
\emph{L'Exercice dans le Christianisme} (1850), annoncèrent le grand règlement
des comptes qu'il allait demander à l'Église. L'idée maîtresse
de ces deux ouvrages est qu'on a oublié, surtout chez les Protestants,
l'idée primitive, le prototype pour le réconciliateur. On a fait du
dogme expiatoire une donnée acquise et on s'est cantonné sans
rancune dans ce monde imparfait. Kierkegaard, tourné vers les
représentants de l'Église, les exhorte à reconnaître quelle distance
sépare la religion actuelle du Nouveau Testament, et dans quelle
apostasie elle est tombée depuis l'aube héroïque du christianisme.
Les chefs d'Église interpellés gardèrent le silence et, même, avancèrent
à l'occasion la prétention d'être en continuité avec « l'époque
des apôtres » : là-dessus Kierkegaard lança une protestation fulgurante,
surtout dans sa brochure belliqueuse \emph{Le Moment} (1855). Mais
% ^ The closing parenthesis of "(1855)" is a battered sort: its upper arm
%   has failed and it prints as a thick baseline hook, which is why the OCR
%   reads it "J". Damaged type, not a different character; set as ")".
il tomba en pleine lutte et mourut. Il expira sans le sacrement, ne
voulant pas le prendre de la main d'un pasteur. Il va plus loin ici
dans sa révolte contre l'Église que ne faisait Pascal, mais il n'a pas
rejeté le sacrement comme tel et, sur ce point, il ne sort pas non plus
des traditions de la vieille Église.

\emph{La Maladie à la mort} est le désespoir qui naît nécessairement
de l'opposition entre la vie spirituelle imposée à l'hommé par le
% ^ PRINTER'S ERROR, as printed: "l'hommé" for "l'homme". The acute on the
%   final e is unmistakable at 300 dpi and again on a 4x enlargement of the
%   word; it is not a speck --- it has the wedge shape of the accent sort.
christianisme et la vie que lui assignent les sentiments et les appétits
naturels. Dans les cas où l'homme n'a pas déjà pris conscience du
dilemme qui l'enserre, il pourra encore jouir de la vie comme le ferait
un enfant ; mais, de loin en loin, l'angoisse s'annoncera mettant fin à
% --- p. 240 ---
sa sécurité. Alors, par faiblesse, par repliement sur lui-même ou
par défi, il pourra tenter d'atténuer sa peur et de rester sourd à
l'exigence chrétienne, mais l'homme n'évitera pas ainsi la présence
invisible du désespoir et ne le vaincra qu'il n'ait au préalable éclaté
et tout envahi. En pratique, la chrétienté s'est habituée --- inconscience,
entraînement romantique des esprits ou simple hypocrisie ---
à réduire le grand conflit entre l'ange et la bête par des concessions
au lieu d'en venir à bout par des victoires.

Dans le livre qu'il donna sous le titre d'\emph{Exercice dans le Christianisme},
Kierkegaard pose en principe que, tant que le Sauveur
n'est pas revenu dans sa gloire, il s'agit pour le chrétien de prendre
pour modèle le Jésus humilié. Nous devons faire abstraction de ce
que le dogme nous transmet sur la nature divine de Jésus et tâcher
de nous rendre contemporains de Jésus de Nazareth, de l'homme
qui se disait être le Fils de Dieu et qui fut en conséquence insulté et
exécuté. Au lieu de cela, on s'est --- tels des écoliers --- procuré par
tricherie la solution du problème, s'évitant la peine de faire le calcul
soi-même. Le parallélisme s'impose avec l'état d'âme de Gethsémani
prescrit par Pascal.

De plus en plus, Kierkegaard en vient à juger que les hommes se
sont rendus coupables d'une grande trahison à l'égard du christianisme :
on s'approprie la douceur et la grâce en oubliant ce qui en
était la condition nécessaire. Il estime que cette subreption est
surtout manifeste dans le protestantisme. Jamais le catholicisme ne
pourrait s'abaisser à une platitude pareille. Et cependant la pratique
des compromissions remonte au catholicisme. L'esprit du christianisme
se trouve faussé le jour où de suivre le Christ à la lettre en
se dérobant à la vie du monde cessa d'être la règle ordinaire. La
différenciation des chrétiens en deux catégories distinctes est un
non-sens. Mais le régime des abus inauguré de la sorte fut repris
par Luther jetant le froc aux orties et fondant une Église nouvelle
où fusionneraient dans un arrangement confortable le christianisme
et le monde. « Luther, Luther, Luther, la responsabilité qui
t'incombe est grande ! » Et Luther était passé par l'école du désespoir.
Mais les générations subséquentes se réclament des actions
d'éclat des ancêtres sans manifester la moindre disposition à en
accomplir elles-mêmes. Le Protestantisme est une doctrine plébéienne
qui tend à niveler, par en bas, le monde de l'esprit. Et pour
une large part c'est la vie de famille qui est la coupable. Par le
% --- p. 241 ---
moyen du baptême des enfants, tous les habitants des pays soi-disant
chrétiens furent enrégimentés dans la chrétienté. « On ose,
s'écrie Kierkegaard (\emph{Le Moment}, VII), on ose faire cela à la face de
Dieu, sous le couvert du baptême chrétien ! Le baptême, l'acte sacré
par lequel le Sauveur du monde fut consacré à l'œuvre de sa vie,
et, après lui, les disciples, des hommes depuis longtemps déjà en
âge de savoir et qui, morts à la vie d'ici-bas, promettaient de vivre
dorénavant comme des sacrifiés dans ce monde du mensonge et de
la méchanceté ! »
% ^ printed without a space before the exclamation mark here
%   ("méchanceté! »"), the line being tightly set; elsewhere on this page
%   the journal's usual space is present ("de Dieu !"). Normalised to the
%   journal's own spacing, which is what the rest of this transcription
%   follows; the compression is justification, not orthography.

La polémique de Kierkegaard eut pour apogée l'énoncé d'une proposition
unique (contrairement aux quatre-vingt-quinze de Luther) :
« Le christianisme du Nouveau Testament n'existe pas ». Et, complétant
la thèse, Kierkegaard mettait en garde contre « ceux qui portent
de longs vêtements » et contre la participation au culte officiel,
où on se moque de Dieu !

L'opposition de Pascal et de Kierkegaard est un fait du plus haut
intérêt dans l'histoire de la religion. Du côté des sociologues, on
a soutenu que la religion est essentiellement un phénomène social.
Et, en fait, l'histoire des religions étudie de préférence des phénomènes
religieux relevant du culte traditionnel. Dans ces conditions,
comment peut-on attribuer une importance du même ordre à des
personnages qui ont été amenés, par l'effort de leur vie et de leur
pensée, à rompre avec la société religieuse à laquelle ils appartenaient,
sans entrer ensuite dans quelque autre société ? Si encore
ils se laissaient rattacher chacun à sa secte : Pascal au jansénisme,
Kierkegaard, éventuellement, au piétisme (comme cela a été fait
dans un excellent ouvrage de sociologie français). Le rattachement
échouerait, d'abord, pour Pascal qui, justement, dans la phase
suprême et décisive de sa vie, eut la douleur de se séparer de ses
alliés de Port-Royal ; et quant à Kierkegaard, il abonde en saillies
ironiques et en sarcasmes à l'adresse des piétistes et ne paraît nullement
enclin à faire cause commune avec eux. Il serait plus plausible
de dire que leur rupture avec la tradition régnante dans
l'Église de leur temps n'eut pour but que d'en relever une plus
ancienne, à leurs yeux la seule légitime. Mais leur appel, leur
retour vers ce qu'ils croyaient l'état de choses originel, en faisait
des solitaires de leur époque, et c'est au nom de la conscience individuelle
qu'ils protestèrent. A quelle instance sociale eussent-ils pu
% ^ "A quelle instance" --- capital A unaccented, as printed.
en appeler ? Pascal a vu, il est vrai, dans un miracle la confirma%
% At the foot of p. 241, below the text and outside it, the printer's
% signature line:  Rev. Méta. --- T. XXX (n° 2, 1923).        16
% (small caps, the "16" at the outer edge). Not part of the article; not set.
% --- p. 242 ---
tion de sa doctrine ; seulement, le miracle en question n'avait pas
reçu d'homologation ecclésiastique. Et Kierkegaard, de son côté,
affirmait expressément que le type de religion qui, seul, avait
pour lui de la valeur, n'existait plus. Dans leurs attitudes
respectives vis-à-vis du sacrement de la dernière heure, un lien
subsiste, trop faible pourtant pour les ramener au contact de
l'Église. --- Au reste, c'est une théorie insoutenable que de s'occuper
à ce point de la tradition et du culte qu'on en vient à négliger
les péripéties antérieures des âmes justiciables de ces forces sociales.
Il faut qu'une des tâches dominantes de l'étude des religions soit
de relever pas à pas les étapes des grandes âmes religieuses. Et
cela d'autant plus que la ligne évolutive de ces individualités supérieures
parcourt la religion en cause jusqu'à ses confins extrêmes,
justement par leur effort même pour en intégrer l'esprit et la vérité
dans leur existence. C'est un tel effort précisément qu'ont réalisé
Pascal et Kierkegaard dans leur développement religieux.

% Section head, centred, on p. 242, with the same setting as the head of
% section I on p. 225: the numeral and the dash in full capitals, then the
% words with the first letter of each in full capitals and the rest in SMALL
% CAPS, and a full stop at the end.  The preposition is set "A", unaccented,
% as the journal's capitals always are.
\section*{IV. --- \textsc{Partis a prendre.}}

Les conjectures sont vaines sur l'orientation qu'aurait pu prendre
l'activité de ces deux grands esprits, si leur vie s'était prolongée.
Mais est-il sans intérêt d'envisager les voies qui pouvaient s'ouvrir
à tout homme ayant suivi d'une adhésion consciente et renforcée
par sa propre expérience les deux protagonistes, jusqu'au point
exact de leur dernière pensée ?

Trois possibilités se présentent.

On pourrait reprendre les principes du christianisme primitif et
s'efforcer d'en faire les fondements d'une vie menée dans l'attente
intense de l'intervention approchante de Dieu, s'abstenant entièrement
de collaborer à la civilisation et aux œuvres courantes d'intérêt
humain. Il existe peut-être, qui sait ? des hommes engagés dans
ce chemin que suivit à sa manière saint François d'Assise.

Un autre parti à prendre serait de commencer par faire la concession,
exigée par Kierkegaard, sur l'écart entre le christianisme
primitif et celui de nos jours. Il faudrait alors faire un pas de plus
dans ce sens et reconnaître qu'au fond le christianisme moderne
est une autre religion que le christianisme primitif. On éviterait
ainsi la contradiction qui consiste à professer un idéal que, pratiquement,
on renie. Dans cette conception on verrait bien dans la
religion un facteur essentiel de la vie spirituelle, mais on admet%
% --- p. 243 ---
trait, à côté, d'autres idéaux qui en diffèrent par leur origine et par
leur orientation et qui, toutefois, méritent de diriger en partie la
conduite et la conception de la vie. Il sera peut-être malaisé de
coordonner dans un même système les différentes étoiles conductrices,
mais, au moins, la sincérité si chère à nos deux penseurs
demeurera intacte.

Pour ceux à qui ces deux chemins paraissent impraticables, un
troisième reste ouvert : ils recommenceront l'effort des Grecs pour
fonder sur la nature et sur l'expérience humaines les directives
d'une vie morale. Ils auront à rechercher l'idéal et le but respectifs
que comportent les principes ainsi établis. En fait, c'est dans ce
sens que se dirigent depuis deux ou trois siècles les tentatives pour
réaliser une morale à base philosophique. J'ai exprimé jadis cette
idée (\emph{La Morale}, X, 4) sous la forme suivante : « Nous pouvons
reprendre la notion platonicienne de la \emph{justice}, en tant qu'harmonie
personnelle, pour exprimer la liaison intime, l'unité supérieure
de l'affirmation de soi et du dévouement. L'antique notion d'harmonie
doit être élargie et approfondie au moyen des expériences
morales qui se trouvent condensées dans les appréciations portées
sur le caractère par le Stoïcisme, le Christianisme, la Renaissance
et l'Humanité moderne. Mais ce changement dans le contenu de la
notion est très possible sans qu'il soit besoin d'en faire éclater le
cadre. Notre morale est la morale grecque, que l'évolution morale
ultérieure, avec ses oscillations multiples, a bien pu modifier et
rectifier, mais qui ne pourra jamais être supplantée, tant qu'il subsistera
une morale humaine, vraiment digne de ce nom. » Le travail
continu sur ces bases ne peut que développer la croyance en son
propre avenir, la conviction qu'il a sa place dans l'existence universelle
et qu'il en représente même un élément indispensable à son
plein développement. Pour l'expression plus vivante d'une telle conviction,
il faut avoir recours aux symbolisations du langage poétique.
Aussi la poésie, riche et profonde, est-elle du nombre des
valeurs dont il faut escompter la permanence et l'épanouissement
dans l'avenir.

En face des deux autres conceptions de la vie que nous avons indiquées,
la position à prendre nous semble être celle-ci : au lieu de renoncer
à tout développement naturel de la vie humaine et à toute culture
autonome, à la vie de famille et à toute activité sociale, artistique ou
scientifique, et au lieu de tenter une combinaison superficielle de
% --- p. 244 ---
dogme oriental et de culture occidentale, on prendra, au contraire,
résolument fait et cause pour la nature humaine et la culture occidentale
et on travaillera courageusement à leur avancement. Les
problèmes à résoudre ne feront pas défaut, ni les déceptions ; des
défaillances se produiront. Il pourra surgir des disputes sur les
chemins à prendre et sur l'établissement des idées directrices, mais
l'orientation générale se laissera toujours reconnaître. D'autres conceptions
de la vie seront mises à contribution ; elles fourniront des
exemples, des aliments utiles. L'élan lyrique que met le croyant
dans son attente du royaume de Dieu, toujours proche, aura une
contre-partie plus ou moins vigoureuse dans toute aspiration active
vers un but élevé. Toute œuvre de grande envergure comporte une
certaine tension. S'il fallait établir une comparaison entre l'influence
qu'a eue le christianisme sur la vie humaine en Europe et
celle exercée par Bouddha en Asie, je dirais, pour reprendre une
formule de ma \emph{Philosophie de la religion}, que Bouddha a adouci
l'Asie et que le Christ a appris à l'Europe l'exaltation infinie de l'effort.
Mais, par-dessus tout, l'esprit de charité prêché par le christianisme
et placé par lui (avec une logique quelquefois défaillante) en
tête des autres vertus, aura sa place indiquée dans une morale
instituée sur des bases purement humaines. Aristote a soutenu
déjà que la justice prend sous sa forme la plus élevée le caractère
de l'amour ; la réciproque est vraie aussi : l'amour n'est vraiment
vertu que si, s'unissant à la sagesse, il devient justice. D'une façon
générale, la conception purement humaine est dans les mêmes rapports
avec le christianisme qu'avec les autres mouvements spirituels :
elle s'assimile ce qu'elle juge compatible avec son aspiration
continue. Autant que les civilisations, les grandes religions universalistes
tirent leurs origines de la nature humaine et des conditions
d'existence de l'homme ; les fruits qu'elles portent, la somme de
noblesse, de grandeur, de beauté que nous leur devons appartient
donc, en dernière instance, à l'humanité tout entière et non pas à
une secte unique, si vaste que soit d'ailleurs son extension.

Quel que soit le chemin où l'on se trouve engagé, on s'inspirera
avec profit de l'enseignement des deux grandes figures qui viennent
de nous retenir. Grâce à elles, le centre de gravité de la vie spirituelle
a passé du monde extérieur des autorités et des dogmes dans
le champ intérieur de la personnalité, de ses expériences et de ses
besoins. Le christianisme, conclut Pascal, est la vraie religion
% --- p. 245 ---
parce qu'il répond seul à l'aspiration la plus profonde de la nature
humaine, parce que seul il concilie les contraires les plus grands de
la vie et qu'il en rend raison. Et pour Kierkegaard, l'idée de Dieu
fut, dans les tempêtes de la vie, une ancre de salut dont la nécessité
ne s'impose pas à celui qui se tient toujours en mer calme. L'essentiel,
pour l'un et l'autre, c'est que chacun passe, au cours de la
vie, par le cycle profitable et pénétrant de l'expérience. Mais d'abord,
et surtout, il faut éviter de se tromper et de tromper les autres en
professant des croyances qui ne se fondent pas sur des expériences
et des aspirations personnelles. « Pour les religions, dit Pascal
(fr. 590 Br.), il faut être sincère ; vrais païens, vrais juifs, vrais
% ^ "590" verified at 4x: the first figure is a 5 and the second a 9. The
%   sandbox OCR reads "580"; the reading is settled by the fragment itself,
%   Brunschvicg 590 being "Pour les religions, il faut être sincère...".
chrétiens. » Kierkegaard, de son côté, disait, pendant sa dernière
campagne contre l'Église établie : « Je n'ose pas m'attribuer le nom
de Chrétien, mais je me range sous le drapeau de la sincérité ; c'est
pour elle que je m'expose ».
% ^ printed thus: the closing guillemet BEFORE the full stop, where the
%   Pascal quotation four lines above closes "chrétiens. »" with the stop
%   inside. The journal's own inconsistency; left standing.

Sur un point, toutefois, les hommes engagés dans le troisième
chemin se séparent nettement de Pascal et de Kierkegaard. Les
deux penseurs se rendent compte que la vie, et particulièrement la
vie religieuse, a besoin de pensées qui en maintiennent et dirigent
l'orientation. Mais ils font une distinction tranchée entre pensées
vitales et pensées scientifiques ou, pour parler le langage de Kierkegaard,
entre connaissance « essentielle » et connaissance « non
essentielle ». Sous la connaissance non essentielle, ils rangent la
science, toutes les sciences. Et pourtant il y a des pensées scientifiques
qui peuvent devenir d'une importance capitale pour les
croyances personnelles. Impossible de marquer ici les limites une
fois pour toutes. Indiquons seulement l'influence sourde et peu
remarquée par le grand public qu'exerce sur les idées religieuses
l'essor admirable qu'ont pris les sciences ; sciences naturelles aussi
bien que sciences morales, pendant les derniers siècles. Et signalons
à ce propos le peu fondé de l'opinion qui veut que la personnalité
et la science s'opposent l'une à l'autre comme s'opposeraient
les intérêts qu'elles recouvrent. La science est issue du travail et des
individus et fille de l'aspiration personnelle ; la joie de connaître
compte parmi les sentiments les plus nobles de l'homme. Les
cadres mêmes conscients ou inconscients de la vie intellectuelle
sont fournis, en dernière analyse, par la personnalité, puisque aussi
bien ils dépendent de la nature et de l'esprit humains. D'un autre
côté, la personnalité se développe par le travail scientifique, elle y
% --- p. 246 ---
gagne une intelligence lucide de la vie et de ses conditions, elle y
apprend à s'imposer des tâches qui dépassent de très haut ses
égoïsmes intéressés. Et il en va de toute autre culture comme de la
science. Sans cesse, l'homme est tantôt supérieur, tantôt inférieur à
telle ou telle culture donnée. Il la domine parce que toute discipline
ne justifie sa valeur qu'en contribuant au développement de
la vie personnelle, et inversement la culture domine l'homme en
imposant des tâches qui subordonnent impérieusement les intérêts
individuels au gain, au progrès durables de la collectivité humaine.

\bigskip
% ^ The journal sets a blank line here --- visibly more than a paragraph
%   break, less than a section break --- before the closing paragraph. The
%   same device as on p. 226. No ornament, no rule, no numeral.

Deux grandes figures, deux solitaires sans continuateurs ni disciples.
Mais l'analogie de leurs idées maîtresses et de leur attitude
respective en face des réalités de leur temps contiennent un enseignement
% ^ printed "contiennent", plural, with the singular subject "l'analogie".
%   Høffding's own French; transcribed as set.
pour toute époque. Ils ont, de leur solitude élevée, discerné
de grands et décisifs problèmes dans le champ de l'histoire
des religions et de la psychologie religieuse, et ils ont laissé leur
vie dans leur énergie à poser ces problèmes d'une inoubliable façon.
D'autant plus inoubliable qu'ils y ont versé non seulement de la
passion mais tant de richesse d'idées et une si dense poésie qu'elles
rangent leurs œuvres parmi les plus précieux trésors humains.

% CLOSE OF THE ARTICLE, p. 246, as printed (read off the image, PDF 264):
% the last paragraph is followed by white, then the author's name alone,
% ranged right (its right edge on the measure), set in capitals and small
% capitals with the umlaut Ö:  HARALD HÖFFDING.  --- H and H full caps,
% the rest small caps, closing full stop. There is NO rule above or below
% it, NO place-and-date line, NO "(A suivre)", NO footnote, and nothing
% else on the page; the remaining two-thirds of the page are blank.
\bigskip

\hfill \textsc{Harald Höffding.}

\end{document}
